\documentclass[../../main.tex]{subfiles}
\begin{document}

{\bf [解答]}

\begin{figure}[htb]
  \centering
  \begin{tikzpicture}[x=1.15cm, y=1.15cm, >=stealth,
    every node/.style={font=\small}]
    % 平地（平面 beta）
    \coordinate (G1) at (1.5, -0.7);
    \coordinate (G2) at (10.8, -0.7);
    \coordinate (G3) at (9.2, 1.8);
    \coordinate (G4) at (-0.1, 1.8);
    \fill[gray!8] (G1) -- (G2) -- (G3) -- (G4) -- cycle;
    \draw[thick, gray!60] (G1) -- (G2) -- (G3) -- (G4) -- cycle;
    \node[gray!70, font=\footnotesize, below right=2pt] at (G4) {平地（平面 $\beta$）};

    % 交線 ell
    \coordinate (B) at (2.80, 0.00);
    \coordinate (C) at (5.00, 0.00);
    \coordinate (A) at (9.71, 0.00);
    \draw[line width=1.3pt, red!75!black] (1.8, 0.0) -- (10.4, 0.0)
    node[right, font=\footnotesize] {交線 $\ell$};

    % 各塔の足
    \coordinate (P0) at (0.74, 0.90);
    \coordinate (Q0) at (4.23, 0.55);
    \coordinate (R0) at (3.58, 0.30);

    % 各塔の先端
    \coordinate (P) at (0.74, 3.96);
    \coordinate (Q) at (4.23, 2.42);
    \coordinate (R) at (3.58, 1.32);

    % 平地上の補助線
    \draw[dotted, gray!60] (Q0) -- (A);
    \draw[dotted, gray!60] (R0) -- (B);
    \draw[dotted, gray!60] (R0) -- (C);

    % 塔の柱
    \draw[line width=1.4pt, gray!85!black] (P0) -- (P)
    node[midway, left=1pt, font=\footnotesize] {$p$};
    \draw[line width=1.4pt, gray!85!black] (Q0) -- (Q)
    node[midway, left=1pt, font=\footnotesize] {$q$};
    \draw[line width=1.4pt, gray!85!black] (R0) -- (R)
    node[midway, left=1pt, font=\footnotesize] {$r$};

    % 先端のつくる平面 alpha
    \fill[blue!15, opacity=0.7] (P) -- (Q) -- (R) -- cycle;
    \draw[thick, blue!85!black] (P) -- (Q) -- (R) -- cycle;
    \node[blue!85!black, font=\footnotesize] at (2.5, 3.0) {平面 $\alpha$};

    % 先端を結ぶ直線の延長
    \draw[dashed, blue!80!black] (Q) -- (A);
    \draw[dashed, blue!80!black] (R) -- (B);
    \draw[dashed, blue!80!black] (R) -- (C);

    % 点の描画
    \node[circle, fill, inner sep=1.2pt] at (P0) {};
    \node[circle, fill, inner sep=1.2pt] at (Q0) {};
    \node[circle, fill, inner sep=1.2pt] at (R0) {};

    \node[circle, fill=blue!85!black, inner sep=1.5pt] at (P) {};
    \node[above=2pt, blue!85!black] at (P) {$P$};
    \node[circle, fill=blue!85!black, inner sep=1.5pt] at (Q) {};
    \node[above right=2pt, blue!85!black] at (Q) {$Q$};
    \node[circle, fill=blue!85!black, inner sep=1.5pt] at (R) {};
    \node[above left=2pt, blue!85!black] at (R) {$R$};

    \node[circle, fill=red!85!black, inner sep=1.6pt] at (A) {};
    \node[below=2pt, red!85!black] at (A) {$A$};
    \node[circle, fill=red!85!black, inner sep=1.6pt] at (B) {};
    \node[below=2pt, red!85!black] at (B) {$B$};
    \node[circle, fill=red!85!black, inner sep=1.6pt] at (C) {};
    \node[below=2pt, red!85!black] at (C) {$C$};
  \end{tikzpicture}
  \caption{3本のテレビ塔の先端を含む平面 $\alpha$ と平地 $\beta$ の交線 $\ell$ 上の3点 $A, B, C$}
  \label{fig:1966_3_overview}
\end{figure}

3本のテレビ塔を $P, Q, R$ としてその高さを $p, q, r\ (p \ge q \ge r)$ とする．
題意より等号が全て成立することはない．
さて，3点 $A, B, C$ は $P, Q$ の先端を結ぶ直線と平地の交点，$Q, R$ の先端を結ぶ直線と平地の交点， $R, P$ の先端を結ぶ直線と平地の交点とおいてよい．

すると，$A, B, C$ はテレビ塔の3先端を含む平面内かつ平地を含む平面内にある．
2平面が交わる時その交線は直線となるから，3点 $A, B, C$ は一直線上にある．
よって題意は示された．

\newpage

\end{document}
